\subsection{One marked system for the chain}
\label{sec:patching-marking}
\hypertarget{source-ordinary-patching-marking}{}

We now put the bounds at the individual points to work on a single system. Keep the
ordinary determinant and the finite chain of comparison arcs fixed. The group \(\Gamma_0\) has a positive subgroup cut out by
the CM sign. The residual characters are defined on the positive
subgroup.
In the CM subgroup write \(\bar\chi_i\) for the residual characters just specified. The subgroup used in the marking is the intersection
\[
 H_{\rm res}=\ker\bar\epsilon\cap\bigcap_i\ker\bar\chi_i
\] It is open.
Polarization permutes the inverse residual characters up to their
common cyclotomic multiplier, so this joint kernel is normal.
The finite tests of Section~\ref{subsec:det-test}
and the simple Frobenius spectra are imposed at all the arcs
before choosing the auxiliary levels. The resulting marked
Frobenius elements \(\sigma_j\) lie in the cyclotomic kernel.

At level \(N\), the surjection \(p_N:\Gamma_N\to\Gamma_0\)
identifies the fixed decomposition maps and the real involution.
It takes the tame generator \(\tau_{j,N}\) to \(1\)
and sends the marked Frobenius to \(\sigma_j\) modulo \(H_{\rm res}\)
and arbitrarily close to \(\sigma_j\) on finite quotients of
the pro-\(\ell\) system. The inertia generator and the norm consequently satisfy 
\[
 p_N(\tau_{j,N})=1,\qquad Q_{j,N}\equiv1\pmod{\ell^N}.
\]
 These requirements on split primes are
obtained by Chebotarev in the finite quotients specifying
the Frobenius approximations, including the cyclotomic quotient
of depth \(N\). The element being approximated is already
trivial in that cyclotomic quotient.

\begin{proposition}[A common marked profinite limit for the chain]
\label{prop:marked-profinite-system}
Fix the comparison arcs and auxiliary prime systems of
Proposition~\ref{prop:marked-uniform-dual-tests}, with each
auxiliary Frobenius converging to its test element in the
common quotient and each norm congruent to \(1\bmod\ell^N\).
After passing to the common residual-cut pro-\(\ell\) quotients,
the point actions, ordinary determinants and representations giving
nonsemisimple quotients at thick points factor through the reduced marked groups. There are a finitely
generated profinite group \(G\) and a surjection
\(p:G\twoheadrightarrow\Gamma_0\) compatible with the fixed
local maps and signs. The distinguished pairs satisfy
\[
 p(\phi_j)=\sigma_j,\qquad
 p(\tau_j)=1,\qquad [\phi_j,\tau_j]=1.
\]
The kernel is the closed normal subgroup generated by the \(\tau_j\).
All these comparisons can be read simultaneously at any fixed finite
depth, with the auxiliary levels allowed to depend on the depth.
\end{proposition}

\begin{proof}
\proofstep{Killing the new tame inertia in the residual cut.}
For a profinite group \(J\) write \(P_\ell(J)\) for the
kernel of its maximal pro-\(\ell\) quotient. Put
\(H_N=p_N^{-1}H_{\rm res}\). The kernel of \(p_N\) is contained
in the closed normal subgroup generated by the
\(\tau_{j,N}\) and \(P_\ell(H_N)\). In the Galois description,
killing the new inertia gives the original ramification
conditions. The additional pro-\(\ell\) quotient of the
pullback of \(H_{\rm res}\) retains just the pro-\(\ell\) tame generator
at each new place.

At residue characteristic \(p\ne\ell\), any image of a pro-\(p\) group in a pro-\(\ell\) quotient is trivial, as its finite quotient orders would have to be powers of both primes:
\[
 p^a=\ell^b\quad\Longrightarrow\quad a=b=0 .
\]
This kills wild inertia and, for the same reason, every prime-to-\(\ell\) primary component of tame inertia. Each new split place and its
conjugate give the same normal generator condition.
The local tame relation is
\(\phi_{j,N}^{-1}\tau_{j,N}\phi_{j,N}=\tau_{j,N}^{Q_{j,N}}\)
in this quotient. Here conjugation by \(\phi_{j,N}^{-1}\) acts on the tame generator by the residue cardinality \(Q_{j,N}\).

\proofstep{The pro-\(\ell\) descent at each depth.}
At every comparison arc the characteristic polynomial on \(H_{\rm res}\)
reduces to \((X-1)^n\). The residual image of \(H_{\rm res}\) in the
finite field matrix group consists of unipotents. If \(u=1+N\) in this finite image, then \(N^n=0\) and in
residue characteristic \(\ell\)
\[
 u^{\ell^a}=1+N^{\ell^a}=1\qquad(\ell^a\ge n).
\]
Every element has \(\ell\)-power order. Cauchy's theorem then makes the
whole image an \(\ell\)-group. Each subsequent congruence kernel in
the curve representation is an \(\ell\)-group. Thus the point
action kills \(P_\ell(H_{\rm res})\), with the same \(H_{\rm res}\) for all points.
The same argument applies to the paired and affine actions at a
finite truncation, using the additive kernels of the coefficient
maps.

For continuous determinants at the ordinary levels,
the split multiplicity-free residual characters give the same
factorization. One may see it on the finite Cayley--Hamilton
quotients: writing \(\bar h\) for the group-algebra image of \(h\in H_{\rm res}\) modulo the maximal ideal, its characteristic polynomial yields 
\[
 P_h(X)\bmod\mathfrak m=(X-1)^n,\qquad (\bar h-1)^n=0.
\]
 At each successively deeper finite quotient the image
of \(H_{\rm res}\) therefore has \(\ell\)-power order. The determinant coefficients and ordered local characters thus factor through the respective quotients
\[
 \Gamma_0\twoheadrightarrow\Gamma_0/P_\ell(H_{\rm res}),\qquad \Gamma_N\twoheadrightarrow\Gamma_N/P_\ell(H_N).
\]
\label{cor:ordinary-patching-marking-1}

\proofstep{The closed matrix algebras after descent.}
This reduction also preserves the thick links. The image of every
element of \(P_\ell(H_{\rm res})\) in the integral closed
Cayley--Hamilton algebra is one at every coefficient depth, hence
one in the separated algebra. Its continuous group map factors
through the quotient group. This gives a map from the closed
Cayley--Hamilton algebra of the quotient group back to the original
one: the determinant and local ordered relations descend, and the
kernel of the map to the finite integral algebra is closed.

In the
opposite direction the quotient homomorphism of groups carries
each defining relation to its counterpart and respects closure.
Writing \(\mathcal E(\Gamma)\) for this closed algebra, the
two constructions give
\[
 \mathcal E(\Gamma_0)\ \longrightarrow\
 \mathcal E(\Gamma_0/P_\ell(H_{\rm res}))\
 \longrightarrow\ \mathcal E(\Gamma_0).
\]
The compositions fix the coefficients and each group image.
These topologically generate the algebras, so both compositions
are the identity. Extending scalars to the field of a nonsemisimple quotient preserves that quotient. The argument
also applies to the marked group with \(H_N\) in place of
\(H_{\rm res}\).

\proofstep{A uniform Frattini bound.}
The bound on the number of generators for the pro-\(\ell\)
groups \(H_N/P_\ell(H_N)\) is uniform in \(N\).
In this Galois setting the fixed field of the residual cut has
fixed degree. Class field theory bounds the \(\mathbb F_\ell\)
dimension of the unramified characters by its fixed class group,
and the restrictions at the fixed finite places contribute a
fixed finite dimension by local class field theory. At each new tame place \(w\), the pro-\(\ell\) tame inertia is procyclic. Its contribution is bounded by 
\[
 \dim_{\mathbb F_\ell}\operatorname{Hom}_{\mathrm{cont}}(I_w,\mathbb F_\ell)\le1.
\]

Each of the at most \(2q\) new places of the CM base has at most the fixed extension degree \(d'\) places above it. Hence
\[
 \#\{\text{new places upstairs}\}\le d'\,\#\{\text{new places downstairs}\}\le 2d'q .
\]
The kernel of restriction to all allowed inertia is unramified.
The same bound with no new places applies to \(H_{\rm res}/P_\ell(H_{\rm res})\).
For a pro-\(\ell\) group \(P\), the continuous characters detect exactly its Frattini quotient: 
\[
 \operatorname{Hom}_{\rm cont}(P,\mathbb F_\ell)=\operatorname{Hom}_{\rm cont}\!\left(P/\overline{P^\ell[P,P]},\mathbb F_\ell\right).
\]
 From now on use the quotients, so \(H_N\) and \(H_{\rm res}\)
themselves are pro-\(\ell\) and finitely generated.

\proofstep{One group marking every auxiliary level.}
Choose a free profinite group \(\mathcal F\) with one generator
for each element of the finite quotient \(\Gamma_0/H_{\rm res}\), a
separate generator for the real involution, and enough additional
generators for the largest of these bounds on \(H_N\). Send
the first generators to representatives in \(\Gamma_N\), the
involution generator to the fixed involution, and the remaining
generators to generators of \(H_N\), padding with identities.
These choices give surjections \(\pi_N:\mathcal F\to\Gamma_N\)
with one common composition into \(\Gamma_0/H_{\rm res}\).
Let \(\mathcal U\) be an ultrafilter extending the tails of the
positive integers.

\proofstep{Compatible finite-quotient limits.}
At an open normal subgroup \(V\subset\mathcal F\), only finitely many subgroups of \(\mathcal F/V\) occur as images. The ultrafilter selects the unique \(K_V\) satisfying
\[
 \{N:(\ker\pi_N)V/V=K_V\}\in\mathcal U .
\]
They are compatible, and every projection between two such
subgroups is onto because it is onto for the corresponding
kernels on their common ultrafilter event. Consequently
\[
 K=\varprojlim_V K_V,\qquad G=\mathcal F/K,\qquad
 G/\operatorname{im}(V)=\mathcal F/(KV).
\]
Compactness and surjectivity of the transition maps give
\(K\bmod V=K_V\); closed normality follows from that of the
\(K_V\).

For every finite quotient \(Q\) of \(\Gamma_0\), the
maps
\[
 \mathcal F\xrightarrow{\pi_N}\Gamma_N \xrightarrow{p_N}\Gamma_0\to Q
\]
have finitely many
possibilities. Their ultrafilter value kills \(K\), since at
the corresponding finite depth the images of the level kernels
coincide with the image of \(K\). The resulting maps \(G\to Q\)
are compatible, so define the continuous map \(p:G\to\Gamma_0\). The map is onto all finite quotients, so it is onto
by compactness. The CM sign, cyclotomic character and real
involution on \(G\) are thereby compatible with the base and
have the required finite quotient comparisons.

\proofstep{Local maps and the marked tame relation.}
The fixed local restriction maps also have simultaneous limits.
For the fixed local group \(D\), let \(f:D\to\Gamma_0\) be its given map. Its subgroup over the residual kernel is open, since
\[
 [D:f^{-1}(H_{\rm res})]=[f(D):f(D)\cap H_{\rm res}]\le[\Gamma_0:H_{\rm res}]<\infty .
\] The maps to all the
\(\Gamma_N\) kill its pro-\(\ell\) kernel since it maps into
\(H_N\). The resulting quotient of the local group is finitely
generated. Its open pro-\(\ell\) subgroup has finite Frattini
rank by local class field theory, and the complementary index is
finite. At each fixed \(V\), equality of the kernel images on an ultrafilter event gives 
\[
 \Gamma_N/\pi_N(V)\simeq\mathcal F/(V\ker\pi_N)=\mathcal F/(VK).
\]
 The finite set of maps from the finitely
generated local quotient to this target has one ultrafilter
value.

Taking all \(V\) defines the local map to \(G\)
with agreement for every element of the local group on a
single event at each fixed depth. Apply the same argument to
the finite tuple of elements \((\phi_{j,N},\tau_{j,N})\)
to obtain \((\phi_j,\tau_j)\) in \(G\). The markings at the limit have the specified Frobenius and inertia values
\[
 p(\phi_j)=\sigma_j,\qquad p(\tau_j)=1 .
\] The \(\tau_j\) lie in \(p^{-1}H_{\rm res}\), a
pro-\(\ell\) group because its finite images come from the
\(H_N\). In a fixed finite quotient choose \(h\) so that the order of the image \(\bar\tau_j\) divides \(\ell^h\). For \(N\ge h\) the finite-depth tame relation
passes to
\[
 \bar\phi_j^{-1}\bar\tau_j\bar\phi_j
 =\bar\tau_j^{\,Q_{j,N}}=\bar\tau_j
 \qquad (Q_{j,N}\equiv1\bmod\ell^h).
\]
The commutator is trivial in every finite quotient and hence in \(G\).

\proofstep{Recovering the base quotient from the limit.}
Finally, every finite quotient of \(G\) killing all the
\(\tau_j\) descends to \(\Gamma_0\). Indeed read its quotient
map on \(\mathcal F\) at a fixed finite depth. There it factors
through \(\Gamma_N\) and kills each marked \(\tau_{j,N}\)
on an ultrafilter event, hence factors through \(p_N\).
The maps from the finitely generated group \(\Gamma_0\)
to this finite quotient form a finite set. Its selected
ultrafilter value therefore gives the factorization through
\(p\).

Put \(N_\tau=\overline{\langle\!\langle\tau_j
\rangle\!\rangle}\). The markings give \(N_\tau\subset\ker p\).
The factorization shows
\[
 \ker p\ \subset\
       \bigcap_{\substack{Q\ {\rm finite}\\G/N_\tau\twoheadrightarrow Q}}
       \ker(G\longrightarrow Q)=N_\tau ,
\]
where the equality holds because \(G/N_\tau\) is profinite.
Thus \(\ker p=N_\tau\). This gives the descent used for
the odd cocycles in Section~\ref{sec:patching-cotangent}.
The construction uses the same free group, limits and local
maps at all chain points. The ultrafilter matching can depend
on truncation depth, and no single finite auxiliary level is
required to match all depths.
\label{cor:ordinary-patching-marking-2}
\end{proof}
